% !TeX root = on-rings-of-integers.tex

% Let \(K\) be a number field, and \(\calO_K\) be its ring of integers.


% \subsection{Hermitian vector bundles}

%     \begin{definition}\label{def:hermitian_coherent_sheaves_on_rings_of_integers}
%         A \emph{hermitian coherent sheaf} \(\overline{\calF}\) on \(\Spec \calO_K\) consists of
%         \begin{itemize}
%             \item a finitely generated \(\calO_K\)-module \(\calF\);
%             \item for each embedding \(\sigma: K \hookrightarrow \bbC\), a Hermitian metric \(\|\cdot\|_\sigma\) on the complex vector space \(\calF_\sigma := \calF \otimes_{\sigma} \bbC\).
%         \end{itemize}
%         \Yang{To be revised.}
%     \end{definition}

%     \begin{definition}\label{def:morphisms_between_hermitian_coherent_sheaves}
%         Let \(\overline{\calF}\) and \(\overline{\calG}\) be hermitian coherent sheaves on \(\Spec \calO_K\).
%         A \emph{morphism} \(f: \overline{\calF} \to \overline{\calG}\) is an \(\calO_K\)-module homomorphism \(f: \calF \to \calG\) such that for each embedding \(\sigma: K \hookrightarrow \bbC\), the induced map \(f_\sigma: \calF_\sigma \to \calG_\sigma\) is a contraction with respect to the Hermitian metrics, i.e.,
%         \[
%             \|f_\sigma(v)\|_\sigma \leq \|v\|_\sigma \quad \text{for all } v \in \calF_\sigma.
%         \]
%         \Yang{To be revised.}
%     \end{definition}

%     \begin{definition}\label{def:hermitian_vector_bundle_on_rings_of_integers}
%         A \emph{hermitian vector bundle} \(\overline{\calE}\) on \(\Spec \calO_K\) consists of
%         \begin{itemize}
%             \item a finitely generated projective \(\calO_K\)-module \(\calE\);
%             \item for each embedding \(\sigma: K \hookrightarrow \bbC\), a Hermitian metric \(\|\cdot\|_\sigma\) on the complex vector space \(\calE_\sigma := \calE \otimes_{\sigma} \bbC\).
%         \end{itemize}
%         \Yang{To be revised.}
%     \end{definition}


% \subsection{Line bundles and divisors}

%     \begin{definition}\label{def:hermitian_line_bundle_on_rings_of_integers}
%         A \emph{hermitian line bundle} \(\overline{\calL}\) on \(\Spec \calO_K\) consists of
%         \begin{itemize}
%             \item a line bundle \(\calL\) on \(\Spec \calO_K\), i.e., a finitely generated projective \(\calO_K\)-module of rank 1;
%             \item for each embedding \(\sigma: K \hookrightarrow \bbC\), a Hermitian metric \(\|\cdot\|_\sigma\) on the complex line \(\calL_\sigma := \calL \otimes_{\sigma} \bbC\).
%         \end{itemize}
%         \Yang{To be revised.}
%     \end{definition}

%     \begin{definition}\label{def:arithmetic_degree_of_hermitian_line_bundles}
%         Let \(\overline{\calL}\) be a hermitian line bundle on \(\Spec \calO_K\). 
%         The \emph{degree} of \(\overline{\calL}\) is defined as
%         \[
%             \deg \overline{\calL} := \log \#(\calL / \calO_K s) - \sum_{\sigma: K \hookrightarrow \bbC} \log \|s\|_\sigma,
%         \]
%         where \(s\) is any non-zero section of \(\calL\).
%         \Yang{To be revised}
%     \end{definition}

%     \begin{definition}\label{def:arithmetic_divisor}
%         An \emph{arithmetic divisor} on \(\Spec \calO_K\) is a pair \((D, g)\), where
%         \begin{itemize}
%             \item \(D\) is a Weil divisor on \(\Spec \calO_K\), i.e., a formal sum of closed points with integer coefficients;
%             \item \(g = (g_\sigma)_{\sigma: K \hookrightarrow \bbC}\) is a collection of Green's functions \(g_\sigma: (\Spec \calO_K)_\sigma(\bbC) \setminus \{D\} \to \bbR\) for each embedding \(\sigma: K \hookrightarrow \bbC\).
%         \end{itemize}
%         \Yang{To be revised}
%     \end{definition}


%     \begin{definition}\label{def:arithmetic_degree_of_arithmetic_divisors}
%         Let \((D, g)\) be an arithmetic divisor on \(\Spec \calO_K\). 
%         The \emph{degree} of \((D, g)\) is defined as
%         \[
%             \deg (D, g) := \sum_{x \in \Spec \calO_K} n_x \log N(x) + \sum_{\sigma: K \hookrightarrow \bbC} g_\sigma(x_\sigma),
%         \]
%         where \(D = \sum_{x} n_x x\) and \(N(x)\) is the norm of the closed point \(x\).
%         \Yang{To be revised}
%     \end{definition}

%     \begin{proposition}\label{prop:arithmetic_divisor_line_bundle_correspondence}
%         There is a one-to-one correspondence between isomorphism classes of hermitian line bundles on \(\Spec \calO_K\) and arithmetic divisors on \(\Spec \calO_K\). 
%         Moreover, this correspondence preserves degrees, i.e., for a hermitian line bundle \(\overline{\calL}\) corresponding to an arithmetic divisor \((D, g)\), we have
%         \[
%             \adeg \overline{\calL} = \adeg (D, g).
%         \]
%         \Yang{To be revised}
        
%     \end{proposition}

%     \begin{proposition}\label{prop:exact_sequence_of_several_class_groups}
%         Let \(K\) be a number field with ring of integers \(\calO_K\).
%         There is an exact sequence
%         \[
%             0 \to \frac{\bigoplus_{\sigma: K \hookrightarrow \bbC} \bbR}{\log |\calO_K^\times|} \to \aCl(\calO_K) \to \Cl(\calO_K) \to 0,
%         \]
%         where \(\aCl(\calO_K)\) is the arithmetic class group of hermitian line bundles on \(\Spec \calO_K\), and \(\Cl(\calO_K)\) is the usual class group of \(\calO_K\).
%         \Yang{To be revised}
%     \end{proposition}


% \subsection{Application: geometry of numbers}


Let \(K\) be a number field with ring of integers \(\calO_K\). 

\subsection{Hermitian line bundles over rings of integers}

    \begin{definition}\label{def:hermitian_line_bundles_over_rings_of_integers}
        A \emph{Hermitian line bundle} \(\overline{\calL}\) over \(\Spec \calO_K\) is a pair \((\calL, (\|\cdot\|_\sigma)_{\sigma: K \to \bbC})\) where
        \begin{itemize}
            \item \(\calL\) is a line bundle over \(\Spec \calO_K\), i.e., a finitely generated projective \(\calO_K\)-module of rank 1.
            \item For each embedding \(\sigma: K \to \bbC\), \(\|\cdot\|_\sigma\) is a hermitian metric on the complex vector space \(\calL_\sigma \coloneqq \calL \otimes_{\calO_K, \sigma} \bbC\).
        \end{itemize}
    \end{definition}

    \begin{example}\label{eg:hermitian_line_bundle_on_Q_sqrt_minus_5}
        Let \(K = \bbQ(\sqrt{-5})\) with ring of integers \(\calO_K = \bbZ[\sqrt{-5}]\).
        The class number of \(K\) is 2, and the ideal \(\frakp = (2, 1 + \sqrt{-5})\) is a non-principal ideal in \(\calO_K\).

        \Yang{}
        
    \end{example}

    \begin{example}\label{eg:hermitian_line_bundle_on_Q_sqrt_10}
        Let \(K = \bbQ(\sqrt{10})\) with ring of integers \(\calO_K = \bbZ[\sqrt{10}]\).
        The class number of \(K\) is 2, and the ideal \(\frakp = (2, \sqrt{10})\) is a non-principal ideal in \(\calO_K\).

        \Yang{To be continued.}
    \end{example}

    \begin{definition}\label{def:arithmetic_degree_of_hermitian_line_bundles}
        Let \(\overline{\calL}\) be a hermitian line bundle on \(\Spec \calO_K\). 
        The \emph{degree} of \(\overline{\calL}\) is defined as
        \[
            \deg \overline{\calL} \coloneqq \log \#(\calL / \calO_K s) - \sum_{\sigma: K \hookrightarrow \bbC} \log \|s\|_\sigma,
        \]
        where \(s\) is any non-zero section of \(\calL\).
        \Yang{To be revised}
    \end{definition}

    \begin{definition}\label{def:arithmetic_divisor}
        An \emph{arithmetic divisor} on \(\Spec \calO_K\) is a pair \((D, g)\), where
        \begin{itemize}
            \item \(D\) is a Weil divisor on \(\Spec \calO_K\), i.e., a formal sum of closed points with integer coefficients;
            \item \(g = (g_\sigma)_{\sigma: K \hookrightarrow \bbC}\) is a collection of Green's functions \(g_\sigma: (\Spec \calO_K)_\sigma(\bbC) \setminus \{D\} \to \bbR\) for each embedding \(\sigma: K \hookrightarrow \bbC\).
        \end{itemize}
        \Yang{To be revised}
    \end{definition}

    \begin{definition}\label{def:arithmetic_degree_of_arithmetic_divisors}
        Let \((D, g)\) be an arithmetic divisor on \(\Spec \calO_K\). 
        The \emph{degree} of \((D, g)\) is defined as
        \[
            \deg (D, g) \coloneqq \sum_{x \in \Spec \calO_K} n_x \log N(x) + \sum_{\sigma: K \hookrightarrow \bbC} g_\sigma(x_\sigma),
        \]
        where \(D = \sum_{x} n_x x\) and \(N(x)\) is the norm of the closed point \(x\).
        \Yang{To be revised}
    \end{definition}

    \begin{proposition}\label{prop:arithmetic_divisor_line_bundle_correspondence}
        There is a one-to-one correspondence between isomorphism classes of hermitian line bundles on \(\Spec \calO_K\) and arithmetic divisors on \(\Spec \calO_K\). 
        Moreover, this correspondence preserves degrees, i.e., for a hermitian line bundle \(\overline{\calL}\) corresponding to an arithmetic divisor \((D, g)\), we have
        \[
            \adeg \overline{\calL} = \adeg (D, g).
        \]
        \Yang{To be revised}
        
    \end{proposition}

    \begin{proposition}\label{prop:exact_sequence_of_several_class_groups}
        Let \(K\) be a number field with ring of integers \(\calO_K\).
        There is an exact sequence
        \[
            0 \to \mu_K \to \calO_K^\times \xrightarrow{\log_K} \bbR^{r} \xrightarrow{l} \aCl(\calO_K) \to \Cl(\calO_K) \to 0,
        \]
        where \(\mu_K\) is the group of roots of unity in \(K\), \(\aCl(\calO_K)\) is the arithmetic class group of \(\calO_K\), and \(\Cl(\calO_K)\) is the usual class group of \(\calO_K\).
        \Yang{To be revised}
    \end{proposition}


\subsection{Geometry of numbers}


    \begin{definition}\label{def:successive_minima_geometry}
        Let \(\overline{\calL}\) be a hermitian line bundle over \(\Spec \calO_K\). 
        The \emph{successive minima} \(\lambda_1(\overline{\calL}), \ldots, \lambda_r(\overline{\calL})\) are defined as
        \[
            \lambda_i(\overline{\calL}) \coloneqq \inf \{ \lambda > 0 : \dim_K K \{ s \in H^0(\Spec \calO_K, \calL) : \|s\|_\sigma \leq \lambda \text{ for all } \sigma \} \geq i \}.
        \]
        \Yang{To be revised}
    \end{definition}


\subsection{Riemann-Roch}

    \begin{definition}\label{def:Euler-Minkowski_characteristic_of_line_bundles}
        Let \(\overline{\calL}\) be a hermitian line bundle over \(\Spec \calO_K\). 
        The \emph{Euler-Minkowski characteristic} \(\chi(\overline{\calL})\) is defined as
        \[
            \chi(\overline{\calL}) \coloneqq \log \frac{\vol(B(\overline{\calL}))}{\vol(\calL_\bbR / \calL)},
        \]
        where \(B(\overline{\calL}) = \{ s \in \calL_\bbR : \|s\|_\sigma \leq 1 \text{ for all } \sigma \}\) and the volumes are computed with respect to the Lebesgue measure.
        \Yang{To be revised}
        
    \end{definition}
